%%________________________________________________________________________
%% LEIM | PROJETO
%% 2026 / 2022 / 2013 / 2012
%% Modelo para relatório
%% v05: alteração DEETC para DEI; outros ajustes
%% v04: alteração ADEETC para DEETC; outros ajustes
%% v03: correção de gralhas
%% v02: inclui anexo sobre utilização sistema controlo de versões
%% v01: original
%% PTS / MAR.2022 / MAI.2013 / 23.MAI.2012 (construído)
%%________________________________________________________________________




%%________________________________________________________________________
\myPrefaceChapter{Resumo}
%%________________________________________________________________________

\textit{Back To School} é um \textit{party game} competitivo multijogador em primeira pessoa para Windows PC desenvolvido em Unity 6.3. Cada ronda é composta por duas fases. A primeira decorre numa sala de aula onde os jogadores respondem a questionários de escolha múltipla gerados dinamicamente por inteligência artificial, comunicam através de voz por proximidade e tentam evitar ser apanhados por um vigia controlado pelo computador. A segunda fase consiste numa corrida de obstáculos até à sala seguinte dentro de um limite de tempo, sendo penalizados os jogadores que não conseguem chegar a tempo.

Este trabalho destaca-se pela integração, numa arquitetura \textit{peer-to-peer}, de três sistemas diferentes, nomeadamente a geração de conteúdo para os quizzes através de uma LLM (\textit{Large Language Model}), NPCs (\textit{Non Player Character}) com sistemas de visão e audição e a comunicação por voz de proximidade. A comunicação por voz representa um dos principais elementos do jogo, uma vez que permite a interação entre jogadores, mas também constitui um risco, pois o som produzido pode alertar o vigia.

A solução foi validada através de testes funcionais e de uma avaliação da experiência de jogo com o questionário GEQ (\textit{Game Experience Questionnaire}) feito com 10 participantes. Os resultados indicam que os participantes compreenderam o objetivo do jogo, valorizando a combinação entre perguntas, interação e corrida de obstáculos, confirmando que o ciclo de jogo proposto é jogável e que a integração de IA generativa acrescenta valor à experiência.





%%________________________________________________________________________
\myPrefaceChapter{Abstract}
%%________________________________________________________________________

\textit{Back To School} is a competitive first-person multiplayer party game for Windows PC developed in Unity 6.3. Each round consists of two phases. The first takes place in a classroom where players answer multiple-choice questions dynamically generated by artificial intelligence, communicate via proximity-based voice chat, and try to avoid being caught by a computer-controlled supervisor. The second phase consists of an obstacle course to the next room within a time limit in which the players who fail to arrive on time are penalized.

This work is distinguished by the integration, within a \textit{peer-to-peer} architecture, of three different systems, including content generation for quizzes using a Large Language Model (LLM), Non-Player Characters (NPCs) with vision and hearing systems, and proximity-based voice communication. Voice communication is one of the game's key elements, as it enables interaction among players, but it also poses a risk, since the sound produced may alert the supervisor.

The solution was validated through functional tests and an evaluation of the gaming experience using the GEQ (Game Experience Questionnaire) with 10 participants. The results indicate that the participants understood the game's objective and appreciated the combination of questions, interaction, and obstacle course, confirming that the proposed game loop is playable and that the integration of generative AI adds value to the experience.
